\section{Deneyler}
\label{sec:setup}

\subsection{Veri kümesi ve bölme protokolü}
HiperKvasir'in 23 sınıflı etiketli alt kümesi kullanılmıştır ($n = 10{.}662$
görüntü). Değerlendirme, veri kümesiyle birlikte sağlanan \emph{resmi 5-katlı}
bölme üzerinden yürütülmüştür. Her katın kendi tutulan (held-out) test bölümü
vardır ve beş kat ayrıktır; bu nedenle birleştirilmiş (havuzlanmış) tahminler
tüm veri kümesini tam olarak bir kez kapsar. Sızıntıyı önlemek için \emph{kat
modelleri asla makro-F1 hesabından önce ortalanmaz}; tüm güven aralıkları, ayrık
kat-dışı (OOF) test tahminlerinden tekrar örnekleyerek makro-F1 dağılımını kestiren
önyükleme (bootstrap) yöntemiyle hesaplanır.

\subsection{Ön işleme}
Görüntüler ImageNet ortalaması/standart sapmasıyla normalize edilir; en kısa
kenar 256'ya ölçeklenip eğitimde rastgele, değerlendirmede merkez $224\times224$
kırpma uygulanır. Yeniden boyutlandırma, \texttt{timm} ViT ağırlıklarıyla
tutarlı olacak şekilde \emph{bikübik} aradeğerleme ile yapılır.

\subsection{Eğitim reçetesi}
İnce ayar reçetesi şu bileşenleri içerir: AdamW
eniyileyici~\cite{loshchilov2019adamw} (omurga öğrenme oranı $5\times10^{-5}$,
başlık $10^{-3}$, ağırlık sönümü $0{,}05$), katman bazında azalan öğrenme oranı
(LLRD, $0{,}7$), 5 epoch ısınma ile kosinüs zamanlayıcı, stokastik derinlik
(drop-path $0{,}05$), etiket yumuşatma $0{,}1$, hafif MixUp ($0{,}2$; küçük
lezyonları silmemek için CutMix kapalı), üstel hareketli ortalama (EMA, $0{,}9998$)
ve dengesizlik için ağırlıklı örnekleyici (sınıf frekansının tersiyle örnekleme
yaparak nadir sınıfları daha sık görünür kılan WeightedRandomSampler). Yığın boyutu
32, en fazla 60 epoch, doğrulama makro-F1 üzerinde erken durdurma. Tek bir model
ailesinin önyargısını azaltmak için ViT-B/BEiT-B'de son üç blok, Swin'de son
aşama açılmıştır. Bu düzenleyici bileşenlerden stokastik derinlik (drop-path),
eğitimde bazı katmanları rastgele atlayarak aşırı uyumu sınırlar; MixUp, örnek
çiftlerini ve etiketlerini doğrusal harmanlayarak karar sınırlarını yumuşatır;
etiket yumuşatma, hedef dağılımı tek-sıcak yerine biraz yumuşatarak aşırı güveni
azaltır; EMA ise model ağırlıklarının üstel hareketli ortalamasını tutarak
değerlendirmede daha kararlı bir model sağlar. Tablo~\ref{tab:hparams} ince ayar
hiperparametrelerini özetler.

\begin{table}[H]
  \centering
  \caption{İnce ayar hiperparametreleri.}
  \label{tab:hparams}
  \begin{tabular}{ll}
    \toprule
    Hiperparametre & Değer \\
    \midrule
    Eniyileyici & AdamW \\
    Omurga öğrenme oranı & $5\times10^{-5}$ \\
    Başlık öğrenme oranı & $10^{-3}$ \\
    Katman bazında azalım (LLRD) & $0{,}7$ \\
    Ağırlık sönümü & $0{,}05$ \\
    Zamanlayıcı & kosinüs, 5 epoch ısınma \\
    Stokastik derinlik (drop-path) & $0{,}05$ \\
    Etiket yumuşatma & $0{,}1$ \\
    MixUp / CutMix & $0{,}2$ / kapalı \\
    EMA sönümü & $0{,}9998$ \\
    Yığın boyutu & 32 \\
    Açılan katmanlar & ViT-B/BEiT-B son 3 blok; Swin son aşama \\
    Epoch (maks.) / erken durdurma & 60 / sabır 8 \\
    Görüntü / normalizasyon & $224{\times}224$, ImageNet, bikübik \\
    Karışık hassasiyet & bfloat16 + TF32 \\
    \bottomrule
  \end{tabular}
\end{table}

\subsection{Hesaplama ve verimlilik}
İnce ayar koşuları Google Colab Pro A100 üzerinde yürütülmüştür. Veri yükleyici
paralelliği, TF32 ve bfloat16 otomatik hassasiyet ile cuDNN otomatik ayarı
etkinleştirildiğinde, A100 verimi temel yapılandırmaya göre yaklaşık
$5{,}3\times$ (tekli) -- $5{,}9\times$ (üçlü) artmıştır; üreme
(reproducibility) yaklaşıktır (seed'ler sabit tutulmuştur). Dondurulmuş eleme
koşuları ise A100 dışında çalıştırılarak hesaplama bütçesi korunmuştur.

\subsection{Ölçütler}
Ana ölçüt \emph{makro-F1}'dir; doğruluk, makro kesinlik/duyarlılık, sınıf bazında
metrikler ve karışıklık matrisi destekleyici olarak raporlanır. Eleme aşaması
yalnızca 0.~katta, seçilen en iyi dört yapılandırma ise tam 5-katlı çapraz
doğrulamada değerlendirilmiştir; bu huni, hesaplama bütçesini korur. Bildirilen
tüm sayılar, her kat için kaydedilen değerlendirme çıktılarından üretilmiştir.
